\section{The question}
\label{sec:question}
%==========================================================================

Propionic acidemia is an autosomal recessive disorder of propionate metabolism. Four amino
acids, odd-chain fatty acids and propionate from intestinal flora converge on
propionyl-CoA, which propionyl-CoA carboxylase (PCC) carboxylates onward into the
tricarboxylic acid cycle. When PCC is defective the intermediate accumulates, spills into
alternative routes, and periodically overwhelms them: the patient decompensates, with
hyperammonaemia and metabolic acidosis. Decompensations are frequent, $49$ of $50$
participants in a natural history cohort reported at least one, at a mean annualised
frequency of $1.0$~\cite{natural_history}, % C08
and they carry the neurological damage that defines the disease's course.

mRNA enzyme replacement supplies PCC transiently~\cite{trial_phase12}. % C16
The clinical decision it creates is a \emph{dosing interval}, and the interval is set by a
quantity nobody measures: the patient's residual enzymatic capacity, written $\theta$
throughout. A patient with more residual capacity needs less frequent replacement.

This note asks a narrow question and answers it. Given a clinical record of the
kind the disease actually generates (a marker series, a dosing diary, a dietary record) \emph{can $\theta$ be identified, and does identifying it decide anything?} The answer
is in four parts, and the second is the one that determines what a protocol should do.

\begin{enumerate}[leftmargin=*, itemsep=3pt]
\item \textbf{The dose train is the identification experiment, and it is already being
run.} The therapy is administered on a fixed biweekly schedule and read as treatment; it is
also a known, exogenous, selective input on exactly the parameter in question, and it
supplies forty-two times the excitation the marker's own excursions do. That is what makes
$\theta$ identifiable: monthly outpatient draws across dose cycles reach $17\%$ over two
years. One scheduling choice destroys it, and it is the choice a protocol would make for
convenience: sampling on dosing days, or at any multiple of the dose interval, costs a
factor of twenty-two, and moving the schedule half a day recovers it
(Sections~\ref{sec:excitation} and~\ref{sec:designsec}).

The obstruction is why this works rather than a finding in its own right. The model is
globally structurally identifiable, so the information is in the system; what a monthly
record loses is the transient. On the quasi-steady locus that remains, the degeneracy has
an exact closed form, a rank-one identity in capacity, affinity and the unrecorded load,
and that identity names the input that closes it. The corollary is a limit on reach
rather than on precision: nothing but the dose train supplies enough excitation, so the
design applies to patients receiving the therapy and to no one else
(Section~\ref{sec:whoto}).
\item \textbf{Identifying $\theta$ does not decide the dosing interval.} The decision is
usually posed as the sign of $\mu = \theta - \thc$ against a critical capacity. That
threshold is not a measured quantity: it is defined through two chosen constants, and
across their defensible range it moves by a factor of $24$ and the sign flips for the same
patient with perfect data (Section~\ref{sec:decision}).
\item \textbf{The threshold framing is itself doubtful.} The only endpoint evidence
available reports a smooth dose--response and no threshold at all. This note records that
as a finding rather than adopting a replacement estimand, for reasons given in
Section~\ref{sec:decision}.
\item \textbf{And capacity stops being recoverable exactly where it matters most.} Above
roughly three to four times the managed marker level, two patients differing tenfold in
residual capacity are indistinguishable, which is the severely affected population
(Section~\ref{sec:scope}).
\end{enumerate}

What this note is not. It is a statement about identifiability, a property of a
model and a design together. It is not a statement about patients, and
Section~\ref{sec:discussion} says plainly where the model's own analysis questions its
validity. The biological background, the disease's clinical course and the argument for
treating $\theta$ as the decision variable are in the supplement, \S S1.

\vspace{0.5em}
\hrule
\vspace{1em}

%==========================================================================
